\clearpage
\fancyhead[L]{\small CARLINI. EFFEMERIDI ASTRONOMICHE.}
\hypertarget{native-571-lower}{}
\begin{center}\small Göttingische gelehrte Anzeigen. No. 48. Pages 475..478. 24 March 1814.\\[3pt]\rule{.72\linewidth}{.3pt}\end{center}

{\rightskip=0pt plus 3em \emph{Effemeridi astronomiche di Milano per l'anno} 1814, \emph{calcolate da} \textsc{Francesco Carlini}. \emph{Dalla stamperia reale. Con appendice} 120 and 140 pages in small quarto.\par}

Since we have already described the practical arrangement of these excellent astronomical ephemerides and the \emph{regular} articles included with them in our reviews of the preceding volumes (see no. 95 from last year [p. 556 of this volume]), we now confine our review to the appendix, which is especially rich in material in this volume too. First we find \textsc{Barnabas Oriani}'s observations of the first comet of 1811. They extend from 29 August 1811 to 21 January 1812. Both the original comparisons made with the equatorial sector and the positions derived from them are given, along with parabolic elements that fit these observations very well and the earlier observations of \textsc{von Zach} less closely. What \textsc{Oriani} says about the outward appearance of this remarkable comet agrees with the observations of other astronomers. He also compares it with the comet of 1744, which in this respect was much like the one of 1811. But this discerning astronomer rightly remarks that the number of larger comets observed with telescopes is still far too small, so it would be too risky to try to base theories on some similarities noticed in the appearance of three or four comets, in the hope of plausibly explaining the origin of the various strange phenomena of their atmospheres and tails. On 24 December \textsc{Oriani} thought he saw for the first time a bright nucleus \emph{in the middle} of the nebulosity in the comet's head. When it could not be made out again in the following days, he had doubts again and thought it likely that the supposed nucleus had been a fixed star of the 9th or 10th magnitude. It seems to us well worth taking another look, for that reason, at the place where the comet stood on 24 December (right ascension $306^\circ\,16'\,21''$, declination $1^\circ\,54'\,6''$ north) (that part of the sky is not visible at this time of year). Next come equally detailed observations of the second comet of 1811 by the same astronomer, along with parabolic elements that do not, however, give very good agreement. As is well known, Mr. \textsc{Nicolai} calculated this comet's orbit as an ellipse, which also fits \textsc{Oriani}'s observations quite well even though they were not used in the calculation. --- Observations of the comet of 1812 by \emph{the same author} (from 1 to 25 September), with parabolic elements. --- Next, a series of papers gives several observed planetary positions by \textsc{Oriani} and \textsc{Santini}, which we can touch on only briefly here. From \textsc{Oriani} come the oppositions of Uranus in 1811, Mars in 1811, Vesta in 1811 and Vesta in 1812, the first three with \textsc{Ramsden}'s mural quadrant and the fourth with the equatorial sector. From \textsc{Santini} in Padua come the oppositions of Juno in 1810 and Uranus in 1810, both with \textsc{Ramsden}'s mural quadrant. --- \SourceLeaf{572}Stellar occultations in 1812 and 1813 and the solar eclipse of 31 January 1813, observed by \textsc{Oriani}, with a detailed calculation of the occultation of $\alpha$ in Taurus on 22 October 1812 to determine the errors of \textsc{Burckhardt}'s new lunar tables. Some readers will be glad to see the necessary rules neatly brought together and explained in practice. For some readers, though, it might have been worth pointing out that in this particular case the result for the error in latitude in the tables could not be very precise. --- A letter from Mr. \textsc{Joseph Piazzi} to \textsc{Oriani}, dated 4 July 1812, contains a few more corrections to his large star catalogue and the smaller one in the Libro Sesto della specola astronomica. It also brings the welcome news that this celebrated observer has already brought a completely new revision of the large catalogue close to completion, which astronomers can hardly wait to see. --- Reflections on astronomical clocks by \textsc{Angelo Cesaris} contain, alongside what is already known, some remarks of no little interest on various modifications of the escapement and other parts tried by Italian craftsmen. One could only wish the remarks were a little more detailed. --- Some further observed and calculated oppositions of the new planets by \textsc{Francesco Carlini}, namely observations of Ceres in 1811 and 1812 and of Pallas in 1811, and calculations of the oppositions of Ceres in 1809, 1811 and 1812, Vesta in 1810, Juno in 1810 and Pallas in 1811. All these results, as well as those mentioned above, are of very great value for planetary theory, since they come from very experienced observers and are based on observations made with excellent instruments. --- The volume closes with the three years of meteorological observations made by \textsc{Cesaris} at the Milan Observatory (1810---1812). We allow ourselves only one remark here. The summer of 1811, so unusually hot in the northern parts of Europe, was far less so in Milan. During the latter half of July the highest thermometer reading was only $25^\circ\,4$, and there were no sudden changes at all. By contrast, in Göttingen, for example, the thermometer in the shade rose to $27^\circ$ on 19 July and even to $29^\circ$ on the 20th, then suddenly fell very low. A discouraging observation for enthusiasts of general theories of the weather!
